\chapter{Berinteraksi dengan API}
\label{chap:apis}

API (Application Programming Interface) adalah jembatan yang menghubungkan kode Anda dengan model AI raksasa di server OpenAI.

\section{Mendapatkan API Key}

\begin{enumerate}
    \item Kunjungi \url{https://platform.openai.com/api-keys}.
    \item Buat akun dan login.
    \item Klik "Create new secret key".
    \item Beri nama kunci, misal "BelajarAI".
    \item \textbf{SALIN DAN SIMPAN SEKARANG JUGA.} Anda tidak akan bisa melihat kunci ini lagi setelah jendela ditutup.
\end{enumerate}

\begin{tipbox}
\textbf{Keamanan API Key:}
Jangan pernah bagikan kunci ini ke orang lain atau upload ke GitHub publik. Kunci ini terhubung ke kartu kredit Anda.
\end{tipbox}

\section{Struktur Request JSON}

Secara teknis, Anda mengirim data dalam format JSON ke server OpenAI.

\begin{codeblock}[title=Contoh Request JSON]
{
  "model": "gpt-4-turbo",
  "messages": [
    {"role": "system", "content": "You are a helpful assistant."},
    {"role": "user", "content": "Hello!"}
  ],
  "temperature": 0.7
}
\end{codeblock}

\section{Membuat Chatbot Sederhana dengan Python}

Mari kita buat chatbot CLI (Command Line Interface) yang bisa ngobrol dengan Anda.

\subsection{1. Simpan API Key}
Buat file bernama `.env` di folder proyek Anda (tanpa nama, hanya ekstensi).

\begin{codeblock}
OPENAI_API_KEY=sk-proj-xxxxxxxxxxxxxxxxxxxxxxxx
\end{codeblock}

\subsection{2. Tulis Kode Chatbot}

Buat file `chatbot.py`.

\begin{codeblock}[language=Python]
import os
from dotenv import load_dotenv
from openai import OpenAI

# 1. Load API Key
load_dotenv()
client = OpenAI(api_key=os.getenv("OPENAI_API_KEY"))

# 2. Fungsi Chat
def chat_dengan_ai():
    print("AI Chatbot (Ketik 'keluar' untuk berhenti)")
    conversation_history = [
        {"role": "system", "content": "Kamu adalah asisten AI yang sopan dan pintar."}
    ]

    while True:
        user_input = input("\nAnda: ")

        if user_input.lower() == "keluar":
            break

        # Tambahkan input user ke history
        conversation_history.append({"role": "user", "content": user_input})

        # Kirim ke OpenAI
        response = client.chat.completions.create(
            model="gpt-3.5-turbo",
            messages=conversation_history,
            temperature=0.7
        )

        # Ambil jawaban AI
        ai_reply = response.choices[0].message.content
        print(f"AI: {ai_reply}")

        # Simpan jawaban AI ke history agar konteks terjaga
        conversation_history.append({"role": "assistant", "content": ai_reply})

if __name__ == "__main__":
    chat_dengan_ai()
\end{codeblock}

\section{Streaming Responses}

Pernah lihat di web ChatGPT, tulisannya muncul satu per satu seperti diketik? Itu namanya "Streaming".
Untuk melakukannya, tambahkan parameter `stream=True` di `client.chat.completions.create`.

\section{Cost Management (Menghitung Token)}

Setiap kali Anda mengirim request, Anda membayar per token.

\begin{itemize}
    \item Input (Prompt): Lebih murah.
    \item Output (Completion): Lebih mahal.
\end{itemize}

Model GPT-4 jauh lebih mahal daripada GPT-3.5. Hati-hati dengan loop `while True` yang tidak terkontrol!

\begin{exercisebox}
\textbf{Tugas Coding:}
Modifikasi skrip `chatbot.py` di atas agar AI berperan sebagai "Guru Bahasa Inggris yang Galak". Jika grammar user salah, AI harus memarahinya sebelum memperbaikinya.
\end{exercisebox}
